\chapter{テキスト生成 — サンプリング戦略とデコード}
\label{ch:generation}

この章では、学習したGPTモデルでテキストを生成するプロセスを実装します。単純なgreedyデコードからTop-K、Top-P、Temperatureサンプリングまで、さまざまな戦略を扱います。

\section{自己回帰生成の原理}

言語モデルの生成は\keyterm{自己回帰（autoregressive)}方式で行われます。すなわち、一度に１つのトークンを生成し、そのトークンを再び入力に追加して次のトークンを生成する処理を繰り返す。

\begin{enumerate}
\item プロンプト$[w\_1, \dots, w\_k]$で始まります。
\item モデルに入れて、次のトークン分布$P(w\_{k+1} | w\_1, \dots, w\_k)$を取得します。
\item 分布から$w\_{k+1}$をサンプリング (または argmax)。
\item シーケンスに追加:$[w\_1, \dots, w\_k, w\_{k+1}]$。
\item 終了条件（<eos>または最大長）まで2$\sim$4繰り返します。
\end{enumerate}

\begin{figure}[htbp]
\centering
\begin{tikzpicture}[
  every node/.style={font=\small\sffamily},
  tok/.style={draw=inkblue, fill=inkblue!8, minimum width=0.7cm, minimum height=0.6cm},
  newtok/.style={draw=inkred, fill=inkred!10, minimum width=0.7cm, minimum height=0.6cm},
  model/.style={draw=inkgray, fill=inkgray!10, rounded corners=2pt, minimum width=2cm, minimum height=0.7cm}
]
  % step 1
  \node[tok] (a1) at (0, 0) {The};
  \node[tok, right=0.05cm of a1] (a2) {cat};
  \node[model, right=0.5cm of a2] (m1) {model};
  \node[newtok, right=0.5cm of m1] (n1) {sat};
  \node[inkgray, anchor=south, font=\scriptsize] at (1.5, 0.5) {step 1};

  % step 2
  \node[tok] (b1) at (0, -1.5) {The};
  \node[tok, right=0.05cm of b1] (b2) {cat};
  \node[tok, right=0.05cm of b2] (b3) {sat};
  \node[model, right=0.5cm of b3] (m2) {model};
  \node[newtok, right=0.5cm of m2] (n2) {on};
  \node[inkgray, anchor=south, font=\scriptsize] at (1.5, -1) {step 2};

  % step 3
  \node[tok] (c1) at (0, -3) {The};
  \node[tok, right=0.05cm of c1] (c2) {cat};
  \node[tok, right=0.05cm of c2] (c3) {sat};
  \node[tok, right=0.05cm of c3] (c4) {on};
  \node[model, right=0.5cm of c4] (m3) {model};
  \node[newtok, right=0.5cm of m3] (n3) {the};
  \node[inkgray, anchor=south, font=\scriptsize] at (1.5, -2.5) {step 3};

  \draw[-{Stealth[length=4pt]}, inkblue] (a2) -- (m1);
  \draw[-{Stealth[length=4pt]}, inkblue] (m1) -- (n1);
  \draw[-{Stealth[length=4pt]}, inkblue] (b3) -- (m2);
  \draw[-{Stealth[length=4pt]}, inkblue] (m2) -- (n2);
  \draw[-{Stealth[length=4pt]}, inkblue] (c4) -- (m3);
  \draw[-{Stealth[length=4pt]}, inkblue] (m3) -- (n3);
\end{tikzpicture}
\caption{自己回帰生成プロセス. ステップごとに前のトークンをモデルに入れ、次のトークンを予測.}
\label{fig:autoregressive}
\end{figure}

\section{Greedy サンプリング}

最も単純な戦略は、各ステップで最も確率の高いトークンを選択することです。

\begin{lstlisting}[language=Python]
class GreedySampler:
    def __call__(self, logits: torch.Tensor) -> torch.Tensor:
        # logits: [batch, vocab_size] -> token_ids: [batch]
        return torch.argmax(logits, dim=-1)
\end{lstlisting}

Greedyは決定論的で高速ですが、毎回同じトークンだけを選択して単調なテキストが出る傾向があります。 「The cat sat on the mat. The cat sat on the mat. The cat... ''同じ繰り返しに陥りやすい。

\begin{warnbox}[Greedyのトラップ：繰り返しループ]
Greedyデコードは「安全な」トークンだけを選択して面白いテキストを作成します。より重大な問題は、繰り返しループに陥るということです。 「私は買い物を買うことができます。買い物をすることができます...」同じ現象。

理由：モデルが「store」の後に「to buy」を高い確率で予測し、「to buy」の後に「some store」を高い確率で予測するループが形成されると、greedyはそこから抜け出せない。
\end{warnbox}

\section{Temperature サンプリング}

分布の「鋭さ」を調整するパラメータを導入します。ロジットを温度$T$で割った後のsoftmax。

\begin{formula}[Temperature Sampling]
\[
P\_T(w) = \frac{\exp(\text{logit}(w) / T)}{\sum\_{w'} \exp(\text{logit}(w') / T)}
\]
\begin{itemize}
\item $T < 1$: 分布が鋭くなる(greedyに近づく)
\item $T > 1$: 分布が平坦になる(よりランダム)
\item $T = 1$: 元の分布
\end{itemize}
\end{formula}

\begin{lstlisting}[language=Python]
class TemperatureSampler:
    def __init__(self, temperature: float = 1.0):
        assert temperature > 0
        self.temperature = temperature

    def __call__(self, logits):
        probs = F.softmax(logits / self.temperature, dim=-1)
        return torch.multinomial(probs, num_samples=1).squeeze(-1)
\end{lstlisting}

\begin{figure}[htbp]
\centering
\begin{tikzpicture}[
  scale=0.8,
  every node/.style={font=\small\sffamily}
]
  % T=0.5（鋭さ）
  \begin{scope}[shift={(0, 0)}]
    \node[inkblue, font=\bfseries\sffamily] at (3, 3.5) {T=0.5（鋭さ）};
    \foreach \i/\h in {0/0.3, 1/0.5, 2/3.0, 3/0.4, 4/0.2, 5/0.1, 6/0.05} {
      \draw[inkblue, fill=inkblue!50] (\i*0.5, 0) rectangle (\i*0.5+0.4, \h);
    }
    \node[inkgray, font=\scriptsize, anchor=north] at (1.75, -0.1) {token};
  \end{scope}

  % T=1.0 (原本)
  \begin{scope}[shift={(5, 0)}]
    \node[inkblue, font=\bfseries\sffamily] at (3, 3.5) {T=1.0 (原本)};
    \foreach \i/\h in {0/0.5, 1/0.8, 2/1.8, 3/1.2, 4/0.7, 5/0.4, 6/0.2} {
      \draw[inkblue, fill=inkblue!50] (\i*0.5, 0) rectangle (\i*0.5+0.4, \h);
    }
  \end{scope}

  % T=2.0（平坦）
  \begin{scope}[shift={(10, 0)}]
    \node[inkblue, font=\bfseries\sffamily] at (3, 3.5) {T=2.0（平坦）};
    \foreach \i/\h in {0/0.8, 1/0.9, 2/1.2, 3/1.0, 4/0.85, 5/0.7, 6/0.6} {
      \draw[inkblue, fill=inkblue!50] (\i*0.5, 0) rectangle (\i*0.5+0.4, \h);
    }
  \end{scope}
\end{tikzpicture}
\caption{Temperatureによる分布の変化. Tが低いと分布が鋭くなります。 greedyに近づく, T高いと平らになり、ランダム性が向上します。.}
\label{fig:temperature}
\end{figure}

\section{Top-K サンプリング}

確率の低いトークンが選択されるのを防ぐために、上位$K$個のトークンのみを考慮してください。

\begin{lstlisting}[language=Python]
class TopKSampler:
    def __init__(self, k: int = 50, temperature: float = 1.0):
        self.k = k
        self.temperature = temperature

    def __call__(self, logits):
        scaled = logits / self.temperature
        top_values, top_indices = torch.topk(scaled, self.k, dim=-1)
        probs = F.softmax(top_values, dim=-1)
        sampled = torch.multinomial(probs, num_samples=1).squeeze(-1)
        return top_indices.gather(-1, sampled.unsqueeze(-1)).squeeze(-1)
\end{lstlisting}

\subsection{Top-Kの直観}

ロジット分布から「尾」（確率が非常に低いトークン）を切り取ります。たとえば、語彙が5万個の場合、上位50個だけを考慮すると99.99\%の確率質量を持つトークンの中からサンプリングするわけだ。

Top-K = 1の場合はgreedyと同じです。 Top-K =$V$（語彙サイズ）の場合は、temperatureサンプリングと同じです。

\section{Top-P (Nucleus) サンプリング}

Top-Kは$K$を固定しますが、Top-Pは分布形式に従って考慮するトークンの数を動的に制御します。累積確率が$P$になるまでのトークンのみを考慮してください。

\begin{lstlisting}[language=Python]
class TopPSampler:
    def __init__(self, p: float = 0.9, temperature: float = 1.0):
        self.p = p
        self.temperature = temperature

    def __call__(self, logits):
        scaled = logits / self.temperature
        probs = F.softmax(scaled, dim=-1)
        sorted_probs, sorted_indices = torch.sort(probs, descending=True, dim=-1)
        cumulative = torch.cumsum(sorted_probs, dim=-1)
        # Pを超える位置からマスキング
        sorted_mask = cumulative - sorted_probs > self.p
        sorted_probs = sorted_probs.masked_fill(sorted_mask, 0.0)
        sorted_probs = sorted_probs / sorted_probs.sum(dim=-1, keepdim=True).clamp(min=1e-10)
        sampled = torch.multinomial(sorted_probs, num_samples=1).squeeze(-1)
        return sorted_indices.gather(-1, sampled.unsqueeze(-1)).squeeze(-1)
\end{lstlisting}

\begin{figure}[htbp]
\centering
\begin{tikzpicture}[
  scale=0.7,
  every node/.style={font=\small\sffamily}
]
  % Top-K可視化(左)
  \begin{scope}[shift={(0, 0)}]
    \node[inkblue, font=\bfseries] at (3, 4.5) {Top-K (K=5)};
    \foreach \i/\h/\sel in {0/2.0/1, 1/3.5/1, 2/1.5/1, 3/0.8/1, 4/2.8/1,
                             5/0.3/0, 6/0.2/0, 7/0.15/0, 8/0.1/0, 9/0.05/0} {
      \pgfmathsetmacro{\col}{ifthenelse(\sel == 1, "inkblue!50", "inkgray!20")}
      \draw[fill=\col, draw=inkblue] (\i*0.5, 0) rectangle (\i*0.5+0.4, \h);
    }
    \draw[inkred, dashed, line width=0.8pt] (2.5, 0) -- (2.5, 4);
    \node[inkred, font=\scriptsize, anchor=west] at (2.6, 3.8) {K=5};
  \end{scope}

  % Top-P可視化(右)
  \begin{scope}[shift={(8, 0)}]
    \node[inkblue, font=\bfseries] at (3, 4.5) {Top-P (P=0.9)};
    \foreach \i/\h/\sel in {0/2.0/1, 1/3.5/1, 2/1.5/1, 3/0.8/1, 4/2.8/1,
                             5/0.3/1, 6/0.2/0, 7/0.15/0, 8/0.1/0, 9/0.05/0} {
      \pgfmathsetmacro{\col}{ifthenelse(\sel == 1, "inkblue!50", "inkgray!20")}
      \draw[fill=\col, draw=inkblue] (\i*0.5, 0) rectangle (\i*0.5+0.4, \h);
    }
    \node[inkred, font=\scriptsize, anchor=west] at (3.0, 3.8) {cumsum};
    \node[inkred, font=\scriptsize, anchor=west] at (3.0, 3.4) {$\approx$ 0.9};
  \end{scope}

  % 軸
  \draw[inkblue, ->] (-0.5, 0) -- (5.5, 0) node[right, font=\scriptsize] {token};
  \draw[inkblue, ->] (-0.5, 0) -- (-0.5, 4) node[above, font=\scriptsize] {prob};
  \draw[inkblue, ->] (5.5, 0) -- (10.5, 0) node[right, font=\scriptsize] {token};
\end{tikzpicture}
\caption{Top-K vs Top-P. Top-Kは固定数です, Top-P累積確率 Pになるまで選択.}
\label{fig:topk-topp}
\end{figure}

\subsection{Top-K vs Top-P}

\begin{center}
\begin{tabularx}{\textwidth}{lXX}
\toprule
 & \textbf{Top-K} & \textbf{Top-P (Nucleus)} \\
\midrule
考慮トークン数 & 固定 ($K$) & 動的 (分布に応じて) \\
分布がシャープである場合＆多すぎるトークンを考慮し、少ないトークンのみを考慮する（自動）\\
分布が平らであるとき＆あまりにも少ない考慮＆多くのトークンを考慮する（自動）\\
適応性 & 低 & 高 \\
\bottomrule
\end{tabularx}
\end{center}

Top-Pは通常より自然な結果を生み出します。モデルが確信する時は少ないトークンだけ、不確実な時は多くのトークンを考慮するから。

\section{サンプリング戦略の比較}

\begin{center}
\begin{tabularx}{\textwidth}{lXXX}
\toprule
\textbf{戦略}&\textbf{多様性}&\textbf{一貫性}&\textbf{ケース}\\
\midrule
Greedy & Low & High & コード生成、数学 \\
Temperature (T=0.7) & 中 & 中 & 一般用途 \\
Top-K (K=50) & 中 & 高 & チャットボット \\
Top-P (P=0.9) & 中 & 高 & チャットボット, 創作 \\
Top-P (P=0.95) & 高 & 中 & 創造的な執筆 \\
\bottomrule
\end{tabularx}
\end{center}

\section{生成関数の実装}

すべてのサンプラーを組み合わせた\code{generate}関数を作成します。

\begin{lstlisting}[language=Python]
@torch.no_grad()
def generate(model, tokenizer, prompt, max_new_tokens=50,
             sampler=None, device="cpu", stop_on_eos=True):
    if sampler is None:
        sampler = GreedySampler()
    model.eval()
    model.to(device)
    ids = tokenizer.encode(prompt, add_bos=True, add_eos=False)
    input_ids = torch.tensor([ids], dtype=torch.long, device=device)
    eos_id = tokenizer.special.eos_id
    context_length = model.config.context_length

    for _ in range(max_new_tokens):
        # コンテキスト長制限
        x = input_ids[:, -context_length:] if input_ids.shape[1] > context_length else input_ids
        logits = model(x)
        next_logits = logits[:, -1, :]  # 最後の場所
        next_id = sampler(next_logits)
        input_ids = torch.cat([input_ids, next_id.unsqueeze(0)], dim=1)
        if stop_on_eos and next_id.item() == eos_id:
            break
    return tokenizer.decode(input_ids[0].tolist())
\end{lstlisting}

\section{コンテキスト長制限}

モデルは学習時に決められた\code{context\_length}までしか処理できません。生成中にシーケンスがこの長さを超えると、最も古いトークンが切り捨てられます。これは近似ですが機能します。 2巻では、この限界を克服するRoPEのextrapolation技術を扱う。

\begin{warnbox}[Sliding Windowの問題]
コンテキストを切り取ると、モデルは「前部」を忘れます。 `` Once upon a time, there was a princess named Aurora. She lived in a castle. One day、[長い説明] ... Auroraは、「Aurora」が前の部分を切り取ると、「Aurora」がプリンセスの名前であるという文脈を失います。

解決策：より長いコンテキストモデルの使用（RoPE extrapolation）、または要約による圧縮。 2巻で扱う。
\end{warnbox}

\section{サンプラーテスト}

\begin{lstlisting}[language=Python]
def test_greedy_picks_max():
    sampler = GreedySampler()
    logits = torch.tensor([[1.0, 3.0, 2.0, 0.5]])
    out = sampler(logits)
    assert out.item() == 1  # 最大値

def test_top_k_limits_choices():
    torch.manual_seed(0)
    logits = torch.tensor([[1.0, 5.0, 4.0, 3.0, 0.1]])
    sampler = TopKSampler(k=2, temperature=1.0)
    for _ in range(20):
        idx = sampler(logits).item()
        assert idx in (1, 2)  # いつも top-2のいずれか

def test_top_p_limits_choices():
    torch.manual_seed(0)
    logits = torch.tensor([[0.0, 10.0, 9.0, 0.0, 0.0]])
    sampler = TopPSampler(p=0.5, temperature=1.0)
    for _ in range(20):
        idx = sampler(logits).item()
        assert idx == 1  # 最も確率の高いトークンのみ
\end{lstlisting}

\section{練習: 各種サンプリング比較}

\code{scripts/tiny\_train.py}デモを実行すると、5つのサンプリング戦略の出力を比較できます。

\begin{lstlisting}[style=bashstyle]
cd make-llm-basic
python scripts/tiny_train.py --epochs 5 --vocab 500
\end{lstlisting}

出力例：
\begin{lstlisting}[basicstyle=\ttfamily\footnotesize, backgroundcolor=\color{codebg}, frame=single]
プロンプト: 'the'
  [Greedy            ] the cat sat on the mat
  [Temperature=0.5   ] the lazy dog sleeps
  [Temperature=1.0   ] the brown fox jumps
  [Top-K=10          ] the quick fox runs
  [Top-P=0.9         ] the warm sun shines
\end{lstlisting}

\section{高度な復号技術(紹介)}

この本では実装しませんが、実際のシステムで使用される高度な技術を紹介します。

\subsection{Beam Search}
各ステップで$K$の候補シーケンスを維持し、最後に最高のスコアのシーケンスを選択します。翻訳、要約などに適していますが、チャットボットには不適合（多様性不足）。

\subsection{Speculative Decoding}
小さな「ドラフト」モデルが最初に$K$個のトークンを生成すると、大きなモデルが並列に検証されます。当たると$K$個を一度に採用、間違えれば最初から。スループット2$\sim$3倍向上可能。

\subsection{Contrastive Search}
繰り返しを避けながら一貫性を保つ最新の技法。以前のトークン表現との類似度でペナルティを与えて様々なトークンを選択。

\section{要約}

\begin{itemize}
\item 自己回帰生成は一度に1つのトークンを生成し、シーケンスを拡張します。
\item Greedy：最も確率の高いトークン。単調さ、繰り返しループの危険。
\item Temperature：分布の鋭さを調整します。 T 低いほど決定論的。
\item Top-K：親K個だけを考慮。シンプルだが非適応。
\item Top-P：累積確率Pまでのトークンのみを考慮します。分布に適応。
\item コンテキスト長を超えると、最も古いトークンを切り捨てます。
\item Speculative decoding、beam searchなど高度な技術も存在。
\end{itemize}

\section{練習問題}

\begin{enumerate}
\item Greedyデコードが繰り返しループに入る理由を分析してください。どのような分布の場合、ループはハンサムですか？
\item Temperatureが0に近づくと、greedyとどう違うのですか？ 0の場合、どのようなことが起こりますか？
\item Top-KとTop-Pを同時に適用するとどうなりますか？ （すなわち、Top-K候補のうちTop-Pフィルタリング）
\item チャットボットにはどのサンプリングが最適ですか？理由を説明してください。
\item Speculative decodingが速い理由は何ですか？
\end{enumerate}

これで1冊のコア実装が終わった。次章では全体をまとめ、2巻で取り上げる上級テーマを紹介する。
